\documentclass[12pt, twoside]{article}
\input{/Users/chris/github/course-files/docs/ib/preamble}

\fancyhead[LE]{\thepage}
\fancyhead[RO]{Markscheme}
\fancyhead[LO]{La Scuola d'Italia / Dr.\ Huson / IB Math HL \\* 28 May 2026}
\fancyfoot[C]{\thepage}

\begin{document}
\subsubsection*{8.1 Classwork: Rational Functions (HL) --- Markscheme}

\medskip

%================================================================
%--- Problem 1 [12 marks]
%================================================================
\noindent\textbf{1.} $\displaystyle f(x)=\dfrac{x^{2}-3x-4}{x-2}$, \ $x\ne 2$.

\medskip

\textbf{(a)} Vertical asymptote: $x=2$. \hfill \textbf{A1}

\medskip

\textbf{(b)} $x$-intercepts: \ $x^{2}-3x-4=0 \Rightarrow (x-4)(x+1)=0$

$(4,0)$ and $(-1,0)$ \hfill \textbf{A1A1}

$y$-intercept: \ $f(0)=\dfrac{-4}{-2}=2$, so $(0,2)$. \hfill \textbf{A1}

\emph{Accept $x=4$, $x=-1$ and $y=2$ in place of the ordered pairs.}

\medskip

\textbf{(c)} Polynomial division of $x^{2}-3x-4$ by $x-2$: \hfill \textbf{M1}

Quotient $x-1$, remainder $-6$. \hfill \textbf{A1}

Hence $\displaystyle f(x)=x-1+\dfrac{-6}{x-2}=x-1-\dfrac{6}{x-2}$ \hfill \textbf{AG}

Oblique asymptote: $y=x-1$. \hfill \textbf{A1}

\emph{M1 may also be awarded for any other valid method, e.g.\ writing
$x^{2}-3x-4=(x-2)(x+c)+r$ and matching coefficients to find $c=-1$, $r=-6$.}

\medskip

\textbf{(d)} \emph{GDC sketch: two branches separated by $x=2$, each approaching the
oblique asymptote $y=x-1$ at the appropriate end; $x$-intercepts at $-1$ and $4$,
$y$-intercept at $2$.}

Correct overall shape with both asymptotes drawn (dashed). \hfill \textbf{A1}

All three intercepts labelled correctly on the sketch. \hfill \textbf{A1}

\medskip

\textbf{(e)} Sign analysis of $\dfrac{(x-4)(x+1)}{x-2}$ on $(-\infty,-1)$,
$(-1,2)$, $(2,4)$, $(4,\infty)$: \hfill \textbf{M1}

$f(x)>0$ on $\,-1<x<2\,$ or $\,x>4$. \hfill \textbf{A1}

\emph{Accept the interval form $(-1,2)\cup(4,\infty)$. \ FT: award A1 if intervals
are read correctly from the candidate's sketch in (d).}

\bigskip\hrule\bigskip

%================================================================
%--- Problem 2 [10 marks]
%================================================================
\noindent\textbf{2.} $\displaystyle f(x)=\dfrac{1}{x^{2}-2x-8}$, \ $x\ne -2,4$.

\medskip

\textbf{(a)} Factor denominator: $x^{2}-2x-8=(x-4)(x+2)$. \hfill \textbf{(M1)}

Vertical asymptotes: $x=-2$ \ \ and \ \ $x=4$. \hfill \textbf{A1A1}

Horizontal asymptote: $y=0$
\ \ (denominator has higher degree than numerator). \hfill \textbf{A1}

\medskip

\textbf{(b)} $f(0)=\dfrac{1}{-8}=-\dfrac{1}{8}$. \hfill \textbf{A1}

\medskip

\textbf{(c)} The denominator $q(x)=x^{2}-2x-8$ has minimum at $x=1$ (vertex of
the parabola): $q'(x)=2x-2=0 \Rightarrow x=1$. \hfill \textbf{M1}

\noindent So $f(x)=\dfrac{1}{q(x)}$ has a stationary point at $x=1$, with

$f(1)=\dfrac{1}{1-2-8}=-\dfrac{1}{9}$. Coordinates: $\left(1,-\tfrac{1}{9}\right)$.
\hfill \textbf{A1A1}

On the middle branch $-2<x<4$, $q(x)<0$ with $q$ achieving its minimum (most negative
value) at $x=1$, so $\dfrac{1}{q(x)}$ achieves its maximum there: local \emph{maximum}.
\hfill \textbf{R1}

\emph{Accept differentiating $f$ directly using the quotient or chain rule,
$f'(x)=-\dfrac{2x-2}{(x^{2}-2x-8)^{2}}$, and solving $f'(x)=0$.}

\medskip

\textbf{(d)} \emph{Sketch shows three branches: left branch ($x<-2$) rising from
$y=0^{+}$ as $x\to-\infty$ up to $+\infty$ as $x\to -2^{-}$; middle branch
($-2<x<4$) an inverted U through $\left(1,-\tfrac{1}{9}\right)$, with
$f\to-\infty$ at both ends; right branch ($x>4$) from $+\infty$ at $x=4^{+}$
down to $y=0^{+}$ as $x\to+\infty$.}

Correct overall shape (three branches; left and right branches above the
$x$-axis, middle branch entirely below). \hfill \textbf{A1}

Asymptotes (dashed), $y$-intercept and local max labelled. \hfill \textbf{A1}

\bigskip\hrule\bigskip

%================================================================
%--- Problem 3 [8 marks]
%================================================================
\noindent\textbf{3.} $\displaystyle f(x)=\dfrac{x-2}{x^{2}-9}$, \ $x\ne\pm 3$.

\medskip

\textbf{(a)} Factor: $x^{2}-9=(x-3)(x+3)$.

Vertical asymptotes: $x=3$ \ and \ $x=-3$. \hfill \textbf{A1A1}

Horizontal asymptote: $y=0$. \hfill \textbf{A1}

\medskip

\textbf{(b)} $x$-intercept: $x-2=0 \Rightarrow (2,0)$. \hfill \textbf{A1}

$y$-intercept: $f(0)=\dfrac{-2}{-9}=\dfrac{2}{9}$, so $\left(0,\tfrac{2}{9}\right)$.
\hfill \textbf{A1}

\medskip

\textbf{(c)} Critical values for the sign chart: $x=-3$, $x=2$, $x=3$.
\ Sign of $\dfrac{(x-2)}{(x-3)(x+3)}$ in each interval: \hfill \textbf{M1}

\smallskip

\begin{center}
\renewcommand{\arraystretch}{1.2}
\begin{tabular}{c|ccccc}
interval & $x<-3$ & $-3<x<2$ & $x=2$ & $2<x<3$ & $x>3$ \\
\hline
sign of $f(x)$ & $-$ & $+$ & $0$ & $-$ & $+$
\end{tabular}
\end{center}

\smallskip

Solution: $-3<x\le 2$ \ \ or \ \ $x>3$. \hfill \textbf{A1A1}

\emph{Accept $(-3,2]\cup(3,\infty)$. \ Note: $x=\pm 3$ are excluded (undefined);
$x=2$ is included since the inequality is non-strict.}

\bigskip\hrule\bigskip

%================================================================
%--- Problem 4 [6 marks]
%================================================================
\noindent\textbf{4.} $\displaystyle g(x)=\dfrac{ax+4}{3-x}$, \ $x\ne 3$.

\medskip

\textbf{(a)} \textbf{Method 1 (swap and equate).}

Let $y=\dfrac{ay+4}{3-y}$ (swap $x$ and $y$): \ \ $x(3-y)=ay+4$ \hfill \textbf{M1}

$3x-4=y(a+x)$ \ $\Rightarrow$ \ $g^{-1}(x)=\dfrac{3x-4}{x+a}$ \hfill \textbf{A1}

Equate $g^{-1}(x)=g(x)$ and cross-multiply: \hfill \textbf{M1}

$(3x-4)(3-x)=(ax+4)(x+a)$, leading on the constant term to
$-12=4a \Rightarrow a=-3$. \hfill \textbf{A1}

\smallskip

\textbf{Method 2 (M\"obius trace condition).} \ For a M\"obius function
$\dfrac{Ax+B}{Cx+D}$, the condition $g^{-1}=g$ (for all $x$ in the domain) is
equivalent to $A+D=0$. Writing $g$ as $\dfrac{ax+4}{-x+3}$ gives $A=a$, $D=3$,
so $a+3=0 \Rightarrow a=-3$. \hfill (\emph{equivalent M1 A1 M1 A1})

\medskip

\textbf{(b)} Line of symmetry: $y=x$. \hfill \textbf{A1}

If $g^{-1}=g$, then $(p,q)$ on the graph of $g$ implies $(q,p)$ is also on the
graph (since $g(p)=q$ gives $g^{-1}(q)=p$, and $g^{-1}=g$ means $g(q)=p$). The
points $(p,q)$ and $(q,p)$ are reflections of one another across the line $y=x$,
hence the graph is symmetric about $y=x$. \hfill \textbf{R1}

\end{document}
